\subsection{Анализ рынка труда}

Для анализа требований рынка труда, предъявляемых к будущей профессии, сформирована \cref{tab:vacancies}. В качестве целевой должности рассматривается инженер-программист АСУ ТП (ПЛК). Отрасли для трудоустройства: машиностроение, приборостроение, инжиниринг, атомная энергетика и др.

\begin{longtblr}[
    caption = {Анализ требований к кандидатам по данным открытых вакансий},
    label = {tab:vacancies},
]{
    colspec = {X[l,t] X[l,t] X[l,t] X[l,t] X[l,t] X[l,t]},
    hlines, vlines,
    row{1} = {bg=gray9, c, m},
    rows = {font=\footnotesize},
    rowhead = 1,
} 
    \textit{Категория требований} & \textit{Вакансия 1} & \textit{Вакансия 2} & \textit{Вакансия 3} & \textit{Вакансия 4} & \textit{Вакансия 5} \\
    \textit{Должность, Компания, Отрасль} & 
    Инженер-программист ПЛК, ИП Шайков В.Н. (Машиностроение) & 
    Инженер АСУ ТП, ООО НПП «Универсал Прибор» (Приборостроение) & 
    Инженер-программист АСУТП, Завод «ДиКом» (Машиностроение) & 
    Сервисный инженер, ООО «Кронштадт» (Инжиниринг) & 
    Инженер-проектировщик систем автоматизации, АО «Гринатом» (Атомная энергетика) \\
    \textit{Знания} & 
    МЭК 61131-3, ТАУ, пром. электроника & 
    МЭК 61131-3, ТАУ, пром. электроника & 
    Электропривод, ПЛК, датчики & 
    Электротехника, электромеханика & 
    Проектирование, локальные сети \\
    \textit{Умения} & 
    Разработка с нуля, чтение схем & 
    Разработка ПО АСУ ТП, ПНР, тестирование & 
    Проектирование, подбор оборудования, сборка, ПНР & 
    Монтаж, ПНР, ремонт, работа с документацией & 
    Разработка проектной документации \\
    \textit{Программное обеспечение} & 
    CODESYS, Easy Builder Pro, MODBUS, OPC & 
    HMI, SCADA, CODESYS, MODBUS, OPC & 
    TIA Portal v16, Sinamics G120 & 
    Не указано & 
    Не указано \\
    \textit{Иностранные языки} & 
    Не требуется & 
    Не требуется & 
    Английский (чтение документации) & 
    Не требуется & 
    Английский (A2) \\
    \textit{Личные особенности} & 
    Готовность к командировкам & 
    Не указано & 
    Не указано & 
    Ответственность, организованность, готовность к командировкам & 
    Аналитический склад ума, ответственность, работа в команде \\
    \textit{Предлагаемая оплата труда, р./мес} & 
    $120000$₽ & 
    $120000$₽ & 
    $180000$₽ & 
    $170000$₽ & 
    Не указана (для расчетов принято $120000$₽) \\
    \textit{Ссылки} & 
    \url{spb.hh.ru/vacancy/132087435} & 
    \url{spb.hh.ru/vacancy/131156471} & 
    \url{spb.hh.ru/vacancy/131596972} & 
    \url{spb.hh.ru/vacancy/131352250} & 
    \url{spb.hh.ru/vacancy/131769673} \\
\end{longtblr}

% На основании проанализированных вакансий определен соответствующий профессиональный стандарт: 40.178 «Специалист по проектированию автоматизированных систем управления технологическими процессами» (утвержден приказом Минтруда РФ № 723н от 12.10.2021). Данный стандарт описывает трудовые функции, связанные с разработкой и внедрением АСУ ТП, что полностью соответствует профилю выбранных должностей.

\subsection{Анализ разрыва и план по его устранению (индивидуальный план развития)}

Сопоставление текущих компетенций (\cref{tab:swot_analysis}) с требованиями работодателей (\cref{tab:vacancies}) выявляет дефицит в следующих областях:
\begin{itemize}
    \item \textit{Твердые умения:} знание стандарта МЭК 61131-3, опыт работы со специализированным ПО (TIA Portal, CODESYS, SCADA), понимание промышленных протоколов (MODBUS, OPC), владение техническим английским языком на уровне не ниже A2-B1.
    \item \textit{Мягкие умения:} навыки управления проектами и работы в команде, адаптивность к командировкам.
\end{itemize}

Для устранения выявленного разрыва сформирован индивидуальный план развития (\cref{tab:idp}).

\begin{longtblr}[
    caption = {Индивидуальный план развития},
    label = {tab:idp},
]{
    colspec = {X[0.8,l,m] X[1.5,l,m] X[1.5,l,m] X[1,l,m] X[1,l,m] X[0.8,l,m]},
    hlines, vlines,
    row{1} = {bg=gray9, c},
    rowhead = 1,
}
    \textit{Зона развития} & \textit{Конкретная цель} & \textit{Мероприятия по достижению} & \textit{Ключевой показатель} & \textit{Ресурс} & \textit{Срок} \\
    Тв. навыки & Изучение языков стандарта МЭК 61131-3 & Прохождение профильного онлайн-курса & Сертификат об окончании & Время ($40$ ч) & $2$ мес. \\
    Тв. навыки & Освоение среды TIA Portal v16 & Выполнение учебного проекта по автоматизации & Готовый проект & Доступ к ПО, время & $3$ мес. \\
    Тв. навыки & Изучение SCADA-систем и HMI & Разработка интерфейса оператора для ПЛК & Рабочий прототип & Доступ к ПО, время & $2$ мес. \\
    Тв. навыки & Освоение протоколов MODBUS и OPC & Настройка обмена данными на лабораторном стенде & Успешная передача пакетов & Лаб. оборудование & $1$ мес. \\
    Тв. навыки & Повышение уровня технического английского & Изучение документации к оборудованию Siemens & Перевод $5$ мануалов & Время ($2$ ч/нед) & $6$ мес. \\
    Тв. навыки & Изучение ГОСТ по оформлению АСУ ТП & Самостоятельное изучение ГОСТ 34 серии & Составленный комплект ТЗ & Доступ к БД ГОСТ & $1$ мес. \\
    Мягк. навыки & Развитие навыков управления проектами & Изучение методологий Agile/Scrum & Успешное тестирование & Время, литература & $2$ мес. \\
    Мягк. навыки & Улучшение навыков работы в команде & Участие в студенческом хакатоне & Сертификат участника & Время ($3$ дня) & $1$ семестр \\
    Мягк. навыки & Тайм-менеджмент в командировках & Разработка системы удаленного контроля задач & Готовый регламент & Литература & $1$ нед. \\
    Мягк. навыки & Развитие лидерских качеств & Выполнение роли Team Lead в студенческом хакатоне & Занятие призовых мест & Время, ответственность & $1$ семестр \\
\end{longtblr}

\subsection{Оценка потенциального дохода}

По данным анализа рынка труда, заработная плата инженера АСУ ТП варьируется от $50000$ до $250000$ руб. в Санкт-Петербурге и от $27659$ до $450000$ руб. по России. Наименьшая предлагаемая заработная плата среди отобранных вакансий составляет $120000$ руб. до вычета налогов.

Расчет чистого потенциального дохода (ЧД) в месяц после вычета НДФЛ ($13\%$):
\begin{equation}
    \text{ЧД} = \text{ЗП} \cdot \frac{100\% - 13\%}{100\%} = 120000 \cdot 0.87 = 104400 \text{ руб.}
\end{equation}

Расчет часовой ставки оплаты труда ($\text{ЧД}_{\text{час}}$) при $5$-дневной рабочей неделе и $8$-часовом рабочем дне:
\begin{equation}
    \text{ЧД}_{\text{час}} = \frac{\text{ЧД}}{22 \cdot 8} = \frac{104400}{176} \approx 593.18 \text{ руб.}
\end{equation}

\subsection{Выбор компании для трудоустройства и составление резюме}
Сформируем образ целевой компании на основе личных предпочтений. Предпочтительное окружение составляют коллеги в возрасте от $22$ до $40$ лет. Рабочая обстановка должна базироваться в офисе с естественным освещением в рамках средней или крупной компании. Целевая деятельность направлена на программирование, проектирование и руководство. Глобальная профессиональная цель заключается в автоматизации процессов. Предпочтителен высокий уровень ответственности, поскольку он в том числе влияет и на уровень дохода. Географическое положение --- Санкт-Петербург или иные мегаполисы России. 

Результаты анализа сведены в \cref{tab:company_profile}.

\begin{longtblr}[
    caption = {Характеристики компании},
    label = {tab:company_profile},
]{
    colspec = {X[1,j,m] X[1,j,m]},
    hlines, vlines,
    rowhead = 1,
}
    \textit{Сфера деятельности компании} & Автоматизация производств, инжиниринг, машиностроение \\
    \textit{Какие мои знания и навыки будут полезны} & Программирование (C, C++, ПЛК), проектирование (САПР), системный анализ \\
    \textit{Каковы люди, которые в ней работают} & Возраст $22-40$ лет, технические специалисты, инженеры \\
    \textit{Масштаб компании} & Средняя или крупная \\
    \textit{Условия труда} & Офис с окнами, наличие оборудованного рабочего места \\
    \textit{Известность компании} & Региональный или федеральный уровень \\
    \textit{Репутация компании} & Надежный интегратор или производитель \\
    \textit{Кто владеет компанией} & Частная или государственная \\
    \textit{Позиционируемые цели компании} & Модернизация промышленности, внедрение инноваций \\
    \textit{Географическое положение компании} & Санкт-Петербург или крупные мегаполисы РФ \\
    \textit{Название компаний, соответствующих критериям} & ООО НПП «Универсал Прибор», Завод «ДиКом» \\
\end{longtblr}

На основании сформированного профиля составлен проект резюме (\cref{tab:resume}).

\begin{longtblr}[
    caption = {Резюме},
    label = {tab:resume},
]{
    colspec = {X[1,l,m] X[2,l,m]},
    hlines, vlines,
    rowhead = 0,
}
    \SetCell[c=2]{c, bg=gray9} \textit{Должность, на которую Вы претендуете} \\
    \SetCell[c=2]{c} Инженер-программист АСУ ТП \\
    \SetCell[c=2]{c, bg=gray9} \textit{Персональные данные} \\
    Фамилия & Рахметов \\
    Имя & Альберт \\
    Отчество & Рафаилович \\
    Контактный телефон & +7 (904) 670-81-55 \\
    Адрес электронной почты & ac7you@gmail.com \\
    Дата рождения & 03.02.2007 \\
    \SetCell[c=2]{c, bg=gray9} \textit{Образование} \\
    Учебное заведение (полное название, год окончания) & Санкт-Петербургский государственный электротехнический университет «ЛЭТИ» им. В. И. Ульянова (Ленина), $2028$ г. \\
    Квалификация и специальность & Бакалавр, «Управление в технических системах» \\
    \SetCell[c=2]{c, bg=gray9} \textit{Опыт работы} \\
    Год и место работы & Выполняемые обязанности \\
    $2026-$наст., ИП Костылев & Оцифровка навигационных справочников судоходства \\
    \SetCell[c=2]{c, bg=gray9} \textit{Профессиональные навыки и умения} \\
    \SetCell[c=2]{l, wd=1\hsize} {Программирование (C, C++, Fortran, MATLAB), 3D-моделирование (Компас-3D), знание электротехники и механики. Навыки оформления технической документации по ГОСТ.} & \\
    \SetCell[c=2]{c, bg=gray9} \textit{Наиболее важные достижения} \\
    \SetCell[c=2]{l, wd=1\hsize} Самостоятельная разработка и внедрение среды формирования типизированных отчетов. Создание информационных веб-ресурсов. \\
    \SetCell[c=2]{c, bg=gray9} \textit{Интересы, увлечения, хобби} \\
    \SetCell[c=2]{l, wd=1\hsize} Модификация Android-устройств (низкоуровневая прошивка, ядро), автоматизация рутинных процессов. \\
    \SetCell[c=2]{c, bg=gray9} \textit{Личные качества} \\
    \SetCell[c=2]{l, wd=1\hsize} Аналитический склад ума, высокая ответственность, склонность к системному планированию, адаптивность к высоким нагрузкам. \\
\end{longtblr}